\item @@M0@@ is the
structural morphism @@M1@@.
\end{enumerate}

\sgaStableTarget{sga6.e08.sec.4}\section*{4. Riemann--Roch theorem: case of a regular closed immersion}

\sgaStableTarget{sga6.e08.para.4-1}\textbf{4.1.} Let @@M2@@ be a regular closed immersion. The Riemann--Roch theorem for @@M3@@ is essentially relation VII 4.4, which makes it possible to compute the Chern classes of @@M4@@ from the @@M5@@. More precisely, if one introduces the @@M6@@-ring @@M7@@ defined in V 5.5, then the homomorphism from @@M8@@ to @@M9@@ obtained from @@M10@@ is a @@M11@@-homomorphism by VII 2.8 and VII 4.3. The Riemann--Roch theorem then reduces to an analogous formula concerning the (additive) homomorphism @@M12@@ defined by @@M13@@, that is, to a formal relation valid for every @@M14@@-ring. We shall not develop this proof here, in order to avoid still more formalism on @@M15@@-rings (cf. @@M16@@ I 1.5), and shall return to the method of [1], which used blow-up schemes (a method already used in the proof of VII 4.3, and which is also that of the end of @@M17@@).

As indicated in 3.9, one may restrict to the case of a regular immersion of constant codimension @@M18@@ satisfying the condition of VII 4.1, the case of a general regular immersion then following from 3.8. In this section we return to the notation of VII 3.1. We thus have the cartesian square:
@@M19@@
